\documentclass[11pt]{article}
\usepackage[margin=0.82in]{geometry}
\usepackage{amsmath,amssymb,mathrsfs,lmodern}
\usepackage[T1]{fontenc}
\usepackage{needspace}
\usepackage{longtable,booktabs,array,calc}
\usepackage[colorlinks=true,linkcolor=blue,urlcolor=blue]{hyperref}
\setlength{\emergencystretch}{3em}
\setlength{\parskip}{0.3em}
\setcounter{secnumdepth}{0}
\providecommand{\tightlist}{\setlength{\itemsep}{0pt}\setlength{\parskip}{0pt}}
\title{Endpoint saddles and the local classification of eigenvalue--torsion critical domains}
\author{Henry Zweiman}
\date{September 16, 2026}
\begin{document}
\maketitle
\Needspace{8\baselineskip}\section{Abstract}\label{abstract}

We study the scale-invariant products of the first Dirichlet eigenvalue
and a positive power of torsional rigidity, normalized by either volume
or perimeter. At each upper local stability threshold identified by
Amato, Nitsch, Trombetti and Villone, we obtain a nonlinear saddle among
smooth strictly convex domains. We also classify every sufficiently
near-ball critical domain, up to similarities, into finitely many
analytic branches indexed by the multiplicities of two eigenvalues of a
trace-free symmetric matrix. Unbalanced branches cross the threshold.
Balanced branches have nonzero quadratic parameter variation; in the
plane the unique nonball branch lies above the threshold. Consequently
the ball is isolated among critical domains at the endpoint despite
being a saddle, and an interval of strict local but nonglobal ball
maximality follows. The argument combines explicit third- and
fourth-order PDE expansions, an analytic boundary equation, and a
finite-dimensional symmetry reduction. A classical transcendence theorem
for rational-order Bessel zeros excludes exceptional dimensions. The
global convex perimeter maximization problem remains unresolved.

\Needspace{8\baselineskip}\section{1. Introduction and main
results}\label{introduction-and-main-results}

For a bounded domain \(\Omega\subset\mathbb R^n\), let
\(\lambda_1(\Omega)\) be the first Dirichlet eigenvalue of \(-\Delta\),
and let \(T(\Omega)=\int_\Omega w\) where \(-\Delta w=1\) in \(\Omega\)
and \(w=0\) on its boundary. Let \(V(\Omega)=|\Omega|\) and
\(P(\Omega)=\mathcal H^{n-1}(\partial\Omega)\). We consider

\[
F_q(\Omega)=\frac{\lambda_1(\Omega)T(\Omega)^q}{V(\Omega)^{\alpha_q}},
\qquad \alpha_q=\frac{(n+2)q-2}{n},
\qquad
G_q(\Omega)=\frac{\lambda_1(\Omega)T(\Omega)^q}{P(\Omega)^{\beta_q}},
\qquad \beta_q=\frac{(n+2)q-2}{n-1}.
\tag{1.1}
\]

Both are invariant under translations and positive dilations. We write
\(B=B_1\), \(\nu=n/2-1\), \(j=j_{\nu,1}\), \(s=j^2\) and \(x=s/n\).
Their upper local stability thresholds are

\[
q_V=\frac{2(x-1)}{n+2},
\qquad q_P=\frac{4(n-1)x-2n+6}{(n+2)(3n-1)}.
\tag{1.2}
\]

These thresholds and their Hessians are established in
\href{https://doi.org/10.1007/s11118-026-10328-2}{Amato-Nitsch-Trombetti-Villone},
Theorems 1.6 and 6.3 and Corollaries 5.14 and 6.7. At each threshold the
Hessian is negative off the degree-two harmonics, after removing
translations and dilations, and vanishes on that finite-dimensional
kernel. Their Remark 5.15 leaves the nonlinear volume endpoint
undecided. For parameters strictly above the thresholds, their results
give strict local ball maximality in the stated smooth topology. We use
this established spectral and stability information; the results below
concern the nonlinear endpoint and the complete nearby critical set.

\Needspace{6\baselineskip}\textbf{Theorem A (endpoint saddles).} In every dimension \(n\ge2\), the
ball is a saddle for \(F_{q_V}\) and \(G_{q_P}\), even among smooth
strictly convex domains. The first assertion holds at exactly fixed
volume and the second at exactly fixed perimeter. In dimensions at least
three the two signs occur along opposite directions of a degree-two
deformation. In the plane they occur along two paths with the same first
velocity and different second-order corrections.

\Needspace{6\baselineskip}\textbf{Theorem B (local classification).} Fix \(n\ge2\), an integer
\(k\ge4\) and \(0<\gamma<1\). For either family, near its parameter
\(q_*\) and near the ball in \(C^{k,\gamma}\), all unrestricted critical
domains, up to similarities, are the ball and the analytic branches

\[
R_{p,t}(\omega)=1+tH_p(\omega)+O(t^2),\qquad
H_p(\omega)=\sum_{i=1}^p\omega_i^2-\frac pn,\qquad
1\le p\le\lfloor n/2\rfloor.
\tag{1.3}
\]

The domains on the \(p\)th branch have \(O(p)\times O(n-p)\) symmetry
and are strictly convex. With \(m_{r,p}=\langle H_p^r\rangle\), their
parameter slopes for \(p<n/2\) are

\[
q_p'(0)=-\frac{3C_*m_{3,p}}{2\,\partial_qe_2\,m_{2,p}}\ne0,
\tag{1.4}
\]

where the explicitly computed cubic coefficients \(C_*\) and the known
Hessian multipliers \(e_2\) are given below. In even dimensions the
balanced branch \(p=n/2\) instead satisfies
\(q_p(t)=q_*+\eta_nt^2+O(t^4)\) with an explicit nonzero \(\eta_n\). The
opposite signs of \(t\) on that balanced branch are congruent. In
dimension two, \(\eta_2>0\) for both families.

\Needspace{6\baselineskip}\textbf{Corollary C (isolated endpoint criticality).} At \(q=q_*\) the
ball is the only critical domain in a sufficiently small neighborhood,
up to similarities. In the plane, for each \(q>q_*\) sufficiently close
to \(q_*\) there is exactly one nearby nonball similarity class, and for
\(q\le q_*\) sufficiently close there is none. In odd dimensions there
are \((n-1)/2\) nonball classes for either sign of sufficiently small
\(q-q_*\). In even dimensions at least four, the \(n/2-1\) unbalanced
classes occur on both sides and the balanced class occurs on the side
determined by \(\eta_n\). A common sign for \(\eta_n\) in every higher
even dimension is not asserted.

\Needspace{6\baselineskip}\textbf{Corollary D (separation of local and global maximality).} For
each family and dimension there is \(\varepsilon_n>0\) such that, for
\(q_*<q<q_*+\varepsilon_n\), the ball is a strict local maximizer but
not a global maximizer, even among smooth strictly convex domains. No
explicit value of \(\varepsilon_n\) is claimed.

The volume branch lies in a different sign regime from the rigidity
theorem of
\href{https://doi.org/10.1007/s00032-024-00405-9}{Celentano-Nitsch-Trombetti},
whose Corollary 1.3 assumes nonpositive product of the exponents in
\(T^a\lambda_1^b\). Here both exponents are positive. The argument does
not contradict that result.

The general implicit-function and symmetry methods are classical. The
background includes Crandall-Rabinowitz\textquotesingle s simple-kernel
bifurcation theory and the literature on spherical symmetry, including
\href{https://doi.org/10.1007/BF00374697}{Chossat-Lauterbach-Melbourne}.
We supply the reduction rather than invoke a new general bifurcation
principle or claim a new general classification of equivariant
equations. The shape-dependent content is the explicit nonlinear
expansion, the verification of the analytic boundary map and its
complement inverse, and the nonvanishing of the coefficients in every
dimension. The latter uses the classical Siegel theorem as stated by
\href{https://doi.org/10.1155/S0161171295000706}{Lorch-Muldoon}: a
nonzero zero of a Bessel function of rational order is transcendental.

The recent preprint of
\href{https://arxiv.org/abs/2609.17083}{Masiello-Paoli-Salerno}, posted
September 15, 2026, concerns saddle shapes for \(k\)-Hessian torsion
under quermassintegral constraints and for curvature integrals. Its
Theorems 1.1-1.4 concern different functionals and a second-order saddle
mechanism. No claim that ball saddles or shape bifurcations are new
phenomena in general is made here.

The perimeter problem at \(q=1\), including the planar conjecture
\(\lambda_1(\Omega)T(\Omega)/P(\Omega)^2\le j_{0,1}^2/(32\pi)\) for
convex \(\Omega\), is outside our conclusion. Local endpoint instability
and critical-domain classification do not settle that global inequality.

\Needspace{8\baselineskip}\section{2. Analytic boundary equation and
Hessian}\label{analytic-boundary-equation-and-hessian}

\subsection{2.1. The analytic boundary
equation}\label{the-analytic-boundary-equation}

Let \(R=1+h\). Write \(u\) for the positive first Dirichlet
eigenfunction normalized by \(\int_{\Omega_h}u^2=1\), \(w\) for the
torsion function, and \(\kappa\) for the sum of the outward principal
curvatures. Put \(E_F=\log F_q\) and \(E_G=\log G_q\). The Hadamard
formulas give their shape densities

\[
\begin{aligned}
\psi_F&=-\frac{u_\nu^2}{\lambda_1}+q\frac{w_\nu^2}{T}-\frac{\alpha_q}{V},\\
\psi_G&=-\frac{u_\nu^2}{\lambda_1}+q\frac{w_\nu^2}{T}-\frac{\beta_q}{P}\kappa.
\end{aligned}
\tag{2.1}
\]

Normal derivatives are outward; only their squares enter. For a radial
variation \(v\), the product of its normal speed and the surface
Jacobian is \(vR^{n-1}\). Thus, with normalized spherical average as
inner product,

\[
D_hE(q,h)[v]=\langle Z(q,h),v\rangle,
\qquad Z=|S^{n-1}|R^{n-1}\psi.
\tag{2.2}
\]

The PDE normalizations in (2.1) matter: the eigenvalue term uses an
\(L^2\)-normalized eigenfunction. First-variation formulas are also
stated in Theorems 2.7-2.8 of Celentano-Nitsch-Trombetti,
\href{https://doi.org/10.1007/s00032-024-00405-9}{final 2024 article}.

Let \(K=O(n-1)\times\{\pm1\}\) act by rotations on the last \(n-1\)
coordinates and reversal of the first. For \(n=2\) this is reflection in
both coordinate axes. Set

\[
\begin{aligned}
X&=\{h\in C^{k,\gamma}(S^{n-1})^K:\langle h\rangle=0\},\\
Y_F&=C^{k-1,\gamma}(S^{n-1})^K\cap\{\langle f\rangle=0\},\\
Y_G&=C^{k-2,\gamma}(S^{n-1})^K\cap\{\langle f\rangle=0\}.
\end{aligned}
\tag{2.3}
\]

Define \(\Psi(q,h)=Z(q,h)-\langle Z(q,h)\rangle\), using the appropriate
\(Y\). These are real-analytic maps on an open neighborhood of
\((q_*,0)\) in \(\mathbb R\times X\).

Here are the analytic details. Choose a fixed smooth radial cutoff
\(\chi\) equal to zero near the origin and to one near \(r=1\). The map
\(r\omega\mapsto r(1+\chi(r)h(\omega))\omega\) is a \(C^{k,\gamma}\)
diffeomorphism from the ball to \(\Omega_h\) when \(h\) is small.
Pulling the Dirichlet equations back to the fixed ball gives
second-order operators whose coefficients are analytic functions of
\(h\), its first and second derivatives, and the inverse of a matrix
close to the identity. They act boundedly from zero-boundary
\(C^{k,\gamma}\) to \(C^{k-2,\gamma}\).

The Dirichlet Laplacian is invertible on these spaces by the classical
Schauder existence and estimate theorem. The torsion solution therefore
depends analytically on \(h\), by the analytic implicit-function
theorem. For the first eigenpair use the equation together with \(L^2\)
normalization, with the pulled-back volume density included. Its
linearization in \((u,\lambda)\) is invertible: projection onto the
simple ground state determines \(\lambda\), the normalization determines
the ground-state component of \(u\), and the Dirichlet resolvent
determines its orthogonal component. The same implicit-function argument
gives analytic dependence. Continuity of the spectrum and positivity
select the first eigenpair. The pulled-back normal derivatives are
analytic in \(C^{k-1,\gamma}\), and the curvature is analytic in
\(C^{k-2,\gamma}\). Integration, logarithms of positive quantities, and
multiplication preserve analyticity. These observations prove the
mapping assertions without assuming a formal shape Taylor series
converges.

The ball satisfies \(\Psi(q,0)=0\) for every nearby \(q\). Further,
\(\Psi(q,h)=0\) is equivalent to unrestricted criticality on this
symmetry class. Indeed \(Z\) is \(K\)-invariant, so \(\Psi=0\) makes it
a constant. Scale invariance gives

\[
0=D_hE(q,h)[1+h]=\langle Z,1+h\rangle=Z,
\tag{2.4}
\]

because \(\langle h\rangle=0\). Thus \(Z=0\) pointwise, and (2.2), or
(2.1) for arbitrary normal speeds, yields every first-variation
equation. Conversely unrestricted criticality implies \(Z=0\). There is
no residual Lagrange multiplier or restriction to symmetric test
variations in this conclusion.

\subsection{2.2. Linearization and the complement
inverse}\label{linearization-and-the-complement-inverse}

For a spherical harmonic of degree \(\ell\ge2\), put

\[
c_\ell=j\frac{J_{\nu+\ell+1}(j)}{J_{\nu+\ell}(j)},
\qquad d_\ell=\ell-c_\ell.
\tag{2.5}
\]

The denominator cannot vanish because the first Dirichlet eigenspace is
radial and simple. The normalized regular radial Helmholtz extension of
a degree-\(\ell\) boundary harmonic has outward derivative \(d_\ell\).

For a single harmonic \(H_\ell\) the coefficient of
\(t^2\langle H_\ell^2\rangle\) in the logarithmic functional along
\(h=tH_\ell\) is

\[
\begin{aligned}
e^F_\ell(q)&=2d_\ell+2(n-1)-(n+2)q(\ell-1),\\
e^G_\ell(q)&=2d_\ell+n-1+\frac{(n+2)q(n+1-2\ell)}2\\
&\hspace{1cm}-\frac{\beta_q}{2}\{\ell(\ell+n-2)+(n-1)(n-2)\}.
\end{aligned}
\tag{2.6}
\]

For completeness, the eigenvalue\textquotesingle s normalized quadratic
coefficient is \(2d_\ell+n-1\), by the second-order compatibility
calculation of Section 3 with degree two replaced by \(\ell\). The
torsion coefficient is \((n+2)((n+1)/2-\ell)\): its first correction is
\(r^\ell H_\ell/n\), its second boundary value is
\((1/2-\ell)H_\ell^2/n\), and the signed shell integral contributes
\(|S|\langle H_\ell^2\rangle/(2n)\). The volume coefficient is
\(n(n-1)/2\). The perimeter coefficient is
\([\ell(\ell+n-2)+(n-1)(n-2)]/2\). Combining them proves (2.6).
Polarization and orthogonality show that \(L=D_h\Psi(q_*,0)\) acts on
each harmonic as multiplication by \(2e_\ell(q_*)\).

The source\textquotesingle s Corollaries 5.14 and 6.7 and their proofs
give

\[
e_2(q_*)=0,\qquad e_\ell(q_*)<0\quad(\ell\ge3).
\tag{2.7}
\]

This spectral information is an established input, not a new endpoint
conclusion. The \(K\)-invariant harmonics of degree two form exactly a
line, spanned by \(H=\omega_1^2-1/n\) for \(n\ge3\), or by
\(H=\cos2\theta\) for \(n=2\). The invariance eliminates degree one.
Thus \(\ker L=\mathbb RH\). Directly from (2.6),

\[
\partial_q e^F_2=-(n+2),\qquad
\partial_q e^G_2=-\frac{(n+2)(3n-1)}{2(n-1)}.
\tag{2.8}
\]

We also need a bounded inverse on the complement in the actual Hölder
spaces, not merely nonzero formal multipliers. Let \(D\) denote the
Laplace Dirichlet-to-Neumann map on the sphere. On mean-zero functions
it is an isomorphism \(C^{k,\gamma}\to C^{k-1,\gamma}\): the inverse is
the harmonic Neumann solution with its constant fixed to have mean-zero
trace. This follows from the Neumann Schauder theorem.

The difference between the Helmholtz and Laplace Dirichlet-to-Neumann
maps is compact as a map from \(C^{k,\gamma}\) to \(C^{k-1,\gamma}\) on
mean-zero traces. Indeed the harmonic extension \(v\) has zero spherical
mean at every radius. The equation

\[
(-\Delta-j^2)z=j^2v,\qquad z|_S=0,
\tag{2.9}
\]

has a unique solution with zero spherical mean, since the only resonant
eigenspace is radial. The Dirichlet Schauder estimate on this complement
bounds \(z\) in \(C^{k+2,\gamma}\) by the boundary trace in
\(C^{k,\gamma}\). The normal derivative of \(z\) lies in
\(C^{k+1,\gamma}\), whose inclusion in \(C^{k-1,\gamma}\) is compact.
The sum \(v+z\) is the Helmholtz extension. Thus the operator with
multipliers \(c_\ell\) is compact at the required orders.

For \(F\), (2.6) now expresses \(L\) as a compact perturbation of

\[
2[2-(n+2)q_V]D:X\longrightarrow Y_F.
\tag{2.10}
\]

For \(G\), it expresses \(L\) as a compact perturbation of

\[
-\beta_{q_P}(-\Delta_S):X\longrightarrow Y_G.
\tag{2.11}
\]

Both leading operators are invertible on mean-zero invariant functions.
In fact \(j^2>2n\) (the standard Bessel-zero bound also used in source
Remark 1.7) gives \(q_V,q_P>2/(n+2)\), so the coefficient in (2.10) is
nonzero and \(\beta_{q_P}>0\). Hence \(L\) is Fredholm of index zero.
Harmonic testing in (2.6) shows its range is orthogonal to \(H\); the
codimension is exactly one by its index and kernel. If

\[
X_\perp=\{h\in X:\langle hH\rangle=0\},\qquad
Y_\perp=\{f\in Y:\langle fH\rangle=0\},
\tag{2.12}
\]

then \(L:X_\perp\to Y_\perp\) is a bounded isomorphism.

\Needspace{8\baselineskip}\section{3. Cubic endpoint expansion in every
dimension}\label{cubic-endpoint-expansion-in-every-dimension}

\subsection{3.1. A cubic expansion along a degree-two
harmonic}\label{a-cubic-expansion-along-a-degree-two-harmonic}

Let \(h\) be a real spherical harmonic of degree two on \(S^{n-1}\),
write \(\langle\cdot\rangle\) for normalized spherical average, and set

\[
\Omega_t=\{r\omega:0\le r<1+t h(\omega)\},\quad
m_2=\langle h^2\rangle,\quad m_3=\langle h^3\rangle.
\tag{3.1}
\]

Here \(\langle h\rangle=0\) and \(-\Delta_S h=2nh\). Put

\[
d=x-n,\qquad a=d+\frac{n-1}{2},\qquad
B=-x^2+\frac{2n+3}{3}x-\frac{n(n-1)}6.
\tag{3.2}
\]

The necessary expansions are

\[
\begin{aligned}
\frac{\lambda_1(\Omega_t)}{j^2}
 &=1+2a m_2t^2+2B m_3t^3+O(t^4),\\
\frac{T(\Omega_t)}{T(B_1)}
 &=1+\frac{(n+2)(n-3)}2m_2t^2
+\frac{(n+2)(n^2-5n+12)}6m_3t^3+O(t^4),\\
\frac{|\Omega_t|}{|B_1|}
 &=1+\frac{n(n-1)}2m_2t^2+\frac{n(n-1)(n-2)}6m_3t^3+O(t^4),\\
\frac{P(\Omega_t)}{P(B_1)}
 &=1+\frac{n^2-n+2}{2}m_2t^2
+\frac{(n-3)(n^2+2)}6m_3t^3+O(t^4).
\end{aligned}
\tag{3.3}
\]

The eigenvalue and torsion remainders will be justified in Section 3.2.
The volume formula follows by expanding \((1+th)^n\). For perimeter,
expand

\[
\frac{P(\Omega_t)}{P(B_1)}
=\left\langle(1+th)^{n-1}
\sqrt{1+\frac{t^2|\nabla_Sh|^2}{(1+th)^2}}\right\rangle.
\tag{3.4}
\]

Use \(\langle|\nabla_Sh|^2\rangle=2nm_2\) and
\(\langle h|\nabla_Sh|^2\rangle=nm_3\). The latter follows by
integrating \(\nabla_S(h^2)\cdot\nabla_Sh=-h^2\Delta_Sh\).

Taking logarithms in (3.3), no cross terms enter before order four. The
quadratic coefficients of \(\log(F_q(\Omega_t)/F_q(B_1))\) and
\(\log(G_q(\Omega_t)/G_q(B_1))\) are, respectively,

\[
[2x-2-(n+2)q]m_2,
\qquad
\left[2x-1+\frac2{n-1}
-\frac{(n+2)(3n-1)}{2(n-1)}q\right]m_2.
\tag{3.5}
\]

They vanish at (1.2). At their respective critical exponents the cubic
coefficients are

\[
\begin{aligned}
C_V(n,x)&=-2x^2+\frac{2(n+8)}3x-\frac83,\\
C_P(n,x)&=-2x^2+
\frac{2(3n^2+22n-9)}{3(3n-1)}x
-\frac{4(5n-3)}{3(3n-1)}.
\end{aligned}
\tag{3.6}
\]

Thus the two logarithmic variations equal \(C_V(n,x)m_3t^3+O(t^4)\) and
\(C_P(n,x)m_3t^3+O(t^4)\).

\subsection{3.2. Derivation of the PDE expansions and their
remainders}\label{derivation-of-the-pde-expansions-and-their-remainders}

\subsubsection{3.2.1. Eigenvalue}\label{eigenvalue}

Let \(u_0(r)=r^{-\nu}J_\nu(jr)\), interpreted smoothly at zero,
\(p=u_0'(1)\ne0\), and \(I=\int_{B_1}u_0^2\). The standard Rellich
identity, obtained by testing the eigenvalue equation against
\(r\partial_ru_0\), gives

\[
p^2|S^{n-1}|=2j^2 I.
\tag{3.7}
\]

Let \(H_2\) be the regular radial solution of the degree-two Helmholtz
equation at \(j^2\), normalized by \(H_2(1)=1\). Then

\[
H_2(r)=\frac{r^{-\nu}J_{\nu+2}(jr)}{J_{\nu+2}(j)},\quad
H_2'(1)=d=\frac{j^2}{n}-n,\quad
H_2''(1)=-(n-1)d-j^2+2n.
\tag{3.8}
\]

The recurrence \(J_{\nu+2}(j)=nJ_{\nu+1}(j)/j\) and the derivative
recurrence prove the middle identity. The denominator is nonzero;
otherwise a degree-two eigenfunction would belong to the first
eigenspace of the ball, contradicting its simplicity. The last identity
is the radial ODE at \(r=1\). The radial equation also gives

\[
u_0''(1)=-(n-1)p,\qquad
u_0'''(1)=[n(n-1)-j^2]p.
\tag{3.9}
\]

We construct a finite expansion \(U_t=u_0+tu_1+t^2u_2+t^3u_3\) and
\(\ell_t=j^2+\lambda_2t^2+\lambda_3t^3\). Let \(L=-\Delta-j^2\). Set

\[
u_1=-pH_2h,
\tag{3.10}
\]

so \(Lu_1=0\) and \(u_1|_S=-ph\). At second order the equations are

\[
Lu_2=\lambda_2u_0,\qquad
u_2|_S=-h\partial_ru_1-\frac12h^2u_0''(1)=pa h^2.
\tag{3.11}
\]

Green\textquotesingle s formula for a function with boundary value \(g\)
gives \(\int_{B_1}u_0 Lu=\int_S pg\). Therefore the Fredholm
compatibility condition for (3.11) is

\[
\lambda_2 I=p^2a|S^{n-1}|m_2,
\qquad \lambda_2=2j^2a m_2.
\tag{3.12}
\]

The degree-two part of \(u_2\) is \(paH_2\Pi_2(h^2)\), where \(\Pi_2\)
is orthogonal projection onto degree-two harmonics. Its remaining parts
have degrees zero and four. Orthogonality implies

\[
\langle h\partial_ru_2\rangle=pa d\,m_3.
\tag{3.13}
\]

This is independent of the arbitrary addition of a multiple of \(u_0\)
to \(u_2\).

At order three choose \(u_3\) to satisfy

\[
Lu_3=\lambda_2u_1+\lambda_3u_0,\qquad
u_3|_S=-h\partial_ru_2-\frac12h^2\partial_r^2u_1
-\frac16h^3u_0'''(1).
\tag{3.14}
\]

Since \(\int u_0u_1=0\), compatibility and (3.7)-(3.13) give

\[
\lambda_3=2j^2\left[-ad+\frac{H_2''(1)}2
-\frac{n(n-1)-j^2}{6}\right]m_3=2j^2B m_3.
\tag{3.15}
\]

This proves the asserted coefficient algebra without assuming that the
formal series is already an eigenfunction.

Here is the remainder argument. The compatibility conditions ensure that
(3.11) and (3.14) have smooth solutions on the ball. All their data
involve finitely many spherical harmonics (degrees at most six);
separation of variables gives regular radial ODE solutions. These extend
smoothly to a fixed ball of radius \(1+\varepsilon\), because the radial
ODE coefficients are smooth at \(r=1\). Thus the finite \(U_t\) is
defined on \(\Omega_t\) and satisfies

\[
(-\Delta-\ell_t)U_t=O(t^4),\qquad
U_t(1+th,\omega)=O(t^4)
\tag{3.16}
\]

in fixed sufficiently strong norms. The first statement follows by
multiplying the finite polynomials in \(t\) and using (3.10)-(3.14); the
second follows from Taylor expansion at \(r=1\) and the boundary
conditions imposed at each order.

Use an angular boundary lift, supported in \(r>(1+th)/2\), of the second
residual. Its \(H^1\) norm is \(O(t^4)\). Subtracting it produces
\(v_t\in H^1_0(\Omega_t)\) with weak eigenvalue residual \(O(t^4)\).
Domain inclusion between balls of radii \(1\pm C|t|\), compactness of
the positive normalized first eigenfunctions, and simplicity of the
disk/ball ground state show that their pairing with \(v_t\) converges to
a nonzero constant. Testing the residual against that first
eigenfunction gives

\[
\lambda_1(\Omega_t)=\ell_t+O(t^4).
\tag{3.17}
\]

The compactness and pairing argument is written explicitly in Section
4.4 and is dimension-independent. This completes the eigenvalue part of
(3.3).

\subsubsection{3.2.2. Torsion}\label{torsion}

Let \(\mathcal H\) be harmonic extension from \(S^{n-1}\) into \(B_1\)
and \(D=\partial_r\mathcal H|_S\); \(D\) acts by multiplication by the
degree on each spherical harmonic. The torsion coefficients are

\[
w_0=\frac{1-r^2}{2n},\qquad
w_1=\frac{r^2h}{n},\qquad
w_2=-\frac{3}{2n}\mathcal H(h^2),
\tag{3.18}
\]

and \(w_3\) is the harmonic extension of

\[
w_3|_S=\frac{3}{2n}hD(h^2)-\frac1n h^3.
\tag{3.19}
\]

These follow by cancelling the boundary values of
\(w_0+tw_1+t^2w_2+t^3w_3\) at \(r=1+th\) order by order. The sum has
Laplacian exactly \(-1\). Orthogonality and \(Dh=2h\) give

\[
\langle hD(h^2)\rangle=2m_3,
\qquad \langle w_3|_S\rangle=\frac2n m_3.
\tag{3.20}
\]

For any harmonic extension, its integral over the ball is
\(|S^{n-1}|/n\) times its boundary average. To include the moving
boundary, integrate the coefficients over \(B_1\) and expand the signed
shell integrals from \(1\) to \(1+th\). At order two the shell
contributes \(|S^{n-1}|m_2/(2n)\). At order three it contributes
\(|S^{n-1}|(n-5)m_3/(6n)\). In detail, the three cubic shell terms,
divided by \(|S^{n-1}|m_3\), are

\[
-\frac{2n-1}{6n}\quad(w_0),\qquad
\frac{n+1}{2n}\quad(tw_1),\qquad
-\frac3{2n}\quad(t^2w_2).
\tag{3.21}
\]

Adding the ball integrals and dividing by
\(T(B_1)=|S^{n-1}|/[n^2(n+2)]\) gives the second line of (3.3). All
harmonic terms are finite harmonic polynomials and extend across the
sphere. The uncancelled boundary value is \(O(t^4)\); the maximum
principle bounds the actual torsion correction by \(O(t^4)\) uniformly,
and hence also in its integral. This supplies the remainder claimed in
(3.3).

\subsection{3.3. Nonvanishing and the higher-dimensional
saddle}\label{nonvanishing-and-the-higher-dimensional-saddle}

For each fixed integer \(n\ge3\), both expressions in (3.6) are nonzero
quadratic polynomials in \(x\) with rational coefficients; each has
leading coefficient \(-2\). The classical theorem of Siegel says that
every nonzero zero of \(J_\nu\) is transcendental for rational \(\nu\).
The applicable statement is given in the introduction of Lorch and
Muldoon, \emph{Transcendentality of zeros of higher derivatives of
functions involving Bessel functions}, International Journal of
Mathematics and Mathematical Sciences 18(3) (1995), 551-560,
\href{https://onlinelibrary.wiley.com/doi/pdf/10.1155/S0161171295000706}{published
PDF},
\href{https://www.maths.tcd.ie/EMIS/journals/HOA/IJMMS/Volume18_3/647167.pdf}{public
journal mirror}.

Since \(\nu=n/2-1\) is rational, \(j\) is transcendental. Consequently
\(x=j^2/n\) is transcendental: an algebraic value of \(x\) would make
\(j\) algebraic as a root of \(z^2-nx\). Thus

\[
C_V(n,x)\ne0,\qquad C_P(n,x)\ne0.
\tag{3.22}
\]

No claim that these coefficients have a common sign across dimensions is
needed.

Choose

\[
h(\omega)=\omega_1^2-\frac1n.
\tag{3.23}
\]

This is a degree-two harmonic because its homogeneous extension is
\(z_1^2-|z|^2/n\). The spherical moments
\(\langle\omega_1^2\rangle=1/n\),
\(\langle\omega_1^4\rangle=3/[n(n+2)]\), and
\(\langle\omega_1^6\rangle=15/[n(n+2)(n+4)]\) give

\[
m_3=\frac{8(n-1)(n-2)}{n^3(n+2)(n+4)}>0.
\tag{3.24}
\]

These moment identities follow, for example, by integrating the
first-coordinate marginal with density proportional to
\((1-z^2)^{(n-3)/2}\) and applying the beta-integral recurrence. From
(3.6), (3.22) and (3.24), each critical logarithmic variation has a
nonzero cubic leading term. Taking \(t\) of opposite signs proves the
saddle property for each functional. The smooth radial hypersurfaces are
strictly convex for small \(|t|\), since their second fundamental forms
converge uniformly to that of the unit sphere.

For exact constraints use \([|B_1|/|\Omega_t|]^{1/n}\Omega_t\) for \(F\)
and \([P(B_1)/P(\Omega_t)]^{1/(n-1)}\Omega_t\) for \(G\). Both scalar
factors tend to 1, preserve strict convexity and preserve the
scale-invariant functional values. This proves the theorem for all
\(n\ge3\).

\Needspace{8\baselineskip}\section{4. The planar volume
endpoint}\label{the-planar-volume-endpoint}

\subsection{4.1. Planar normalization and the two
paths}\label{planar-normalization-and-the-two-paths}

Write \(B\) for the unit disk, \(j=j_{0,1}\), \(s=j^2\), and

\[
q_* = \frac{s-2}{4},\qquad
F_q(\Omega)=\frac{\lambda_1(\Omega)T(\Omega)^q}{|\Omega|^{2q-1}}.
\tag{4.1}
\]

Here \(-\Delta w=1\) with zero Dirichlet data and \(T=\int w\). This
normalization is important: the exponent in the area denominator is
\(2q-1\).

\textbf{Planar endpoint assertion.} The disk is a saddle for \(F_{q_*}\)
among smooth strictly convex planar domains, including under an exact
area constraint. More precisely, set

\[
\Omega_{t,b}=\{(r,\theta):0\le r<1+t\cos(2\theta)+bt^2\cos(4\theta)\},
\qquad b\in\{0,1\}.
\tag{4.2}
\]

For all sufficiently small nonzero \(t\),

\[
F_{q_*}(\Omega_{t,0})>F_{q_*}(B),\qquad
F_{q_*}(\Omega_{t,1})<F_{q_*}(B).
\tag{4.3}
\]

Both paths have the same first normal velocity \(\cos(2\theta)\). Their
values differ first at fourth order. Multiplication by
\((\pi/|\Omega_{t,b}|)^{1/2}\) imposes exact area \(\pi\) and preserves
the inequalities by scale invariance. These scaled paths still have the
same first normal velocity. In particular, a nonpositive Hessian does
not imply a local maximum at this endpoint.

The domain boundaries are positive radial graphs and converge in every
fixed \(C^k\) norm to the circle. The numerator of their signed
curvature is \(R^2+2(R')^2-RR''\), which converges uniformly to 1. Thus
these are strictly convex smooth domains for small \(|t|\). The unscaled
radial graph perturbations also have zero mean and zero first harmonics,
as in the source\textquotesingle s local class.

\subsection{4.2. Exact expansions and the signs that
suffice}\label{exact-expansions-and-the-signs-that-suffice}

For each fixed real \(b\), the expansions used below are

\[
\begin{aligned}
\frac{|\Omega_{t,b}|}{\pi}&=1+\frac{t^2}{2}+\frac{b^2t^4}{2},\\
\frac{T(\Omega_{t,b})}{\pi/8}
 &=1-t^2+\left(\frac{15}{8}+3b-5b^2\right)t^4+O(t^5),\\
\sqrt{\lambda_1(\Omega_{t,b})}
 &=j+k_2t^2+k_4t^4+O(t^5),\qquad
k_2=\frac{j(s-3)}4,\\
k_4&=\frac{j K_b(s)}{256(s-6)},
\end{aligned}
\tag{4.4}
\]

where

\[
\begin{aligned}
K_b(s)={}&16b^2s^2-576b^2s+2688b^2
-24bs^3+440bs^2-1792bs+384b\\
&-7s^4+110s^3-703s^2+2072s-1920.
\end{aligned}
\tag{4.5}
\]

All remainders are as \(t\to0\) for fixed \(b\); no uniform-in-\(b\)
assertion is needed. Sections 4.3 and 4.4 prove the error estimates and
specify a finite recurrence for checking every coefficient.

Taking logarithms in (4.4), the quadratic coefficient in
\(\log(F_q(\Omega_{t,b})/F_q(B))\) is

\[
\frac{s-3}{2}-q-\frac{2q-1}{2}=\frac{s-2-4q}{2}.
\tag{4.6}
\]

It vanishes at \(q=q_*\). The quartic coefficient there is

\[
\begin{aligned}
128(6-s)C_b(s)={}&176b^2s^2-1024b^2s+24bs^3\\
&-536bs^2+2560bs-1536b\\
&+7s^4-102s^3+555s^2-1280s+768.
\end{aligned}
\tag{4.7}
\]

Consequently

\[
\log\frac{F_{q_*}(\Omega_{t,b})}{F_{q_*}(B)}
= C_b(s)t^4+O(t^5).
\tag{4.8}
\]

In particular, with

\[
\begin{aligned}
P_0(x)&=7x^4-102x^3+555x^2-1280x+768,\\
P_1(x)&=7x^4-78x^3+195x^2+256x-768,
\end{aligned}
\tag{4.9}
\]

we have \(C_0(s)=P_0(s)/(128(6-s))\) and \(C_1(s)=P_1(s)/(128(6-s))\).

Here is an exact sign proof without numerical Bessel evaluation. The
Bessel power series and alternating remainder bounds give

\[
J_0(\sqrt{23/4})\ge\sum_{k=0}^{5}\frac{(-23/16)^k}{(k!)^2}
=\frac{17997619}{5033164800}>0,
\tag{4.10}
\]

and

\[
J_0(\sqrt{29/5})\le\sum_{k=0}^{4}\frac{(-29/20)^k}{(k!)^2}
=-\frac{42533}{30720000}<0.
\tag{4.11}
\]

The terms after the first decrease in magnitude at these arguments,
which is sufficient for the stated tail bounds. Moreover for
\(0<z\le\sqrt{29/5}\), setting \(v=z^2/4\), the series for \(J_1\) has
decreasing alternating terms and yields

\[
J_1(z)\ge\frac z2\left(1-\frac v2\right)>0.
\tag{4.12}
\]

Thus \(J_0'=-J_1<0\) on this interval. Its first positive zero satisfies

\[
\frac{23}{4}<s<\frac{29}{5}<6.
\tag{4.13}
\]

For \(y=x-23/4\) one computes

\[
\begin{aligned}
P_0(23/4+y)&=7y^4+59y^3+\frac{1473}{8}y^2+\frac{4935}{16}y+\frac{4719}{256},\\
P_1(23/4+y)&=7y^4+83y^3+\frac{1905}{8}y^2+\frac{1359}{16}y-\frac{6513}{256}.
\end{aligned}
\tag{4.14}
\]

The first is positive for \(y\ge0\). The second is increasing there, and
its value at \(x=29/5\) is \(-12868/625<0\). Equations (4.13)-(4.14)
prove \(C_0(s)>0>C_1(s)\). Formula (4.8) therefore proves (4.3).

\subsection{4.3. Torsion expansion with a controlled
remainder}\label{torsion-expansion-with-a-controlled-remainder}

Let \(R=1+t\cos2\theta+bt^2\cos4\theta\) and use \(\rho\) as an interior
radial coordinate. Seek

\[
W_t(\rho,\theta)=\frac{1-\rho^2}{4}
+c_0(t)+\sum_{n\in\{2,4,6,8\}}c_n(t)\rho^n\cos(n\theta).
\tag{4.15}
\]

This function satisfies \(-\Delta W_t=1\) exactly. The following
coefficients cancel its boundary values through degree four in \(t\):

\[
\begin{aligned}
c_0(t)&=-\frac38t^2-\frac{14b^2-16b-3}{16}t^4,\\
c_2(t)&=\frac12t-\frac{10b-3}{8}t^3,\\
c_4(t)&=\frac{4b-3}{8}t^2+\frac{3(8b-3)}8t^4,\\
c_6(t)&=-\frac{5(2b-1)}8t^3,\\
c_8(t)&=-\frac{7(2b^2-8b+3)}{16}t^4.
\end{aligned}
\tag{4.16}
\]

These identities can be checked by substituting \(\rho=R\) in (4.15),
expanding powers of \(R\), and using
\(2\cos a\cos b=\cos(a+b)+\cos(a-b)\). Since only finite polynomials are
involved, the uncancelled boundary value is \(O(t^5)\) uniformly in
\(\theta\). The actual torsion function \(w_t-W_t\) is harmonic with
boundary data \(-W_t\). The maximum principle gives
\(\|w_t-W_t\|_\infty=O(t^5)\). The areas are uniformly bounded, so their
integrals differ by \(O(t^5)\).

Writing \(\langle f\rangle=(2\pi)^{-1}\int_0^{2\pi}f\,d\theta\), direct
integration of (4.15) gives

\[
\frac{1}{2\pi}\int_{\Omega_{t,b}}W_t
=\left\langle\frac{R^2}{8}-\frac{R^4}{16}
+\sum_{n\in\{0,2,4,6,8\}}\frac{c_n(t)}{n+2}R^{n+2}\cos(n\theta)\right\rangle.
\tag{4.17}
\]

Inserting (4.16) and keeping terms through degree four proves the
torsion line of (4.4). The area identity follows directly from
\(|\Omega|=\pi\langle R^2\rangle\).

\subsection{4.4. Eigenvalue expansion with a controlled
remainder}\label{eigenvalue-expansion-with-a-controlled-remainder}

Set \(k(t)=j+k_2t^2+k_4t^4\) and consider the finite Bessel sum

\[
\begin{aligned}
U_t(\rho,\theta)={}&J_0(k(t)\rho)
+(A_{21}t+A_{23}t^3)J_2(k(t)\rho)\cos2\theta\\
&+(A_{42}t^2+A_{44}t^4)J_4(k(t)\rho)\cos4\theta\\
&+A_{63}t^3J_6(k(t)\rho)\cos6\theta
+A_{84}t^4J_8(k(t)\rho)\cos8\theta.
\end{aligned}
\tag{4.18}
\]

It is smooth at the origin and satisfies \(-\Delta U_t=k(t)^2U_t\)
exactly. Put

\[
\begin{aligned}
A_{21}&=s/2,\\
A_{42}&=-\frac{s^2(4b-s+3)}{32(s-6)},\\
A_{23}&=-\frac{s(20bs^2-288bs+960b+7s^3-89s^2+416s-768)}{128(s-6)}.
\end{aligned}
\tag{4.19}
\]

The remaining coefficients \(A_{63},A_{44},A_{84}\) are uniquely
specified by cancelling the degree-three mode 6 and degree-four modes 4
and 8 in \(U_t(R,\theta)\). Their explicit expanded rational expressions
are unnecessary for the constant equation. Existence here does not
assume a solution of an infinite system: each coefficient first enters
its indicated harmonic multiplied by \(J_n(j)\ne0\), and there are just
these three scalar equations.

For completeness, nonvanishing follows from the Bessel series: for every
integer \(n\ge1\) and \(0<z\le\sqrt{29/5}\),

\[
J_n(z)\ge \frac{(z/2)^n}{n!}
\left(1-\frac{z^2}{4(n+1)}\right)>0.
\tag{4.20}
\]

The alternating terms decrease, since their successive magnitude ratio
is \((z^2/4)/((k+1)(n+k+1))<1\).

Here is a reproducible finite coefficient calculation. Normalize
\(B_n=J_n(j)/J_1(j)\), with

\[
B_0=0,\quad B_1=1,\quad B_{n+1}=\frac{2n}{j}B_n-B_{n-1},\quad
J_n^{(d)}(j)=2^{-d}\sum_{a=0}^d(-1)^a\binom da J_{n-d+2a}(j),
\tag{4.21}
\]

using \(J_{-n}=(-1)^nJ_n\). Expand the argument increment as

\[
k(t)R-j=j t c+(k_2+jbd)t^2+k_2c t^3+(k_4+k_2bd)t^4+O(t^5),
\quad c=\cos2\theta,\ d=\cos4\theta.
\tag{4.22}
\]

At order one, mode 2 yields \(A_{21}\). At order two, modes 0 and 4
yield \(k_2,A_{42}\). At order three, modes 2 and 6 yield
\(A_{23},A_{63}\). At order four, modes 0, 4 and 8 yield
\(k_4,A_{44},A_{84}\). No other modes occur at these orders. In
particular, the mean equation at order two is

\[
-J_1(j)k_2+\frac{jJ_1(j)}4+\frac{A_{21}jJ_2'(j)}2=0,
\tag{4.23}
\]

which gives (4.4). The mean equation at order four is

\[
\begin{aligned}
0={}&-J_1(j)k_4
+\frac{J_0''(j)}2\left(k_2^2+\frac{j^2b^2}{2}+jk_2\right)
+J_0'''(j)\left(\frac{j^2k_2}{4}+\frac{j^3b}{8}\right)
+\frac{j^4J_0''''(j)}{64}\\
&+A_{21}\left(\frac{k_2J_2'(j)}2
+J_2''(j)\left(\frac{jk_2}{2}+\frac{j^2b}{4}\right)
+\frac{j^3J_2'''(j)}{16}\right)
+\frac{A_{23}jJ_2'(j)}2\\
&+A_{42}\left(\frac{jbJ_4'(j)}2+\frac{j^2J_4''(j)}8\right).
\end{aligned}
\tag{4.24}
\]

It uses \(\langle c^2\rangle=1/2\), \(\langle c^4\rangle=3/8\), and
\(\langle c^2d\rangle=1/4\). Substitution of (4.19) and (4.21) gives
exactly (4.5). This displays the scalar calculation determining the
decisive fourth-order coefficient, rather than relying on a numerical
eigensolver.

After all these choices, \(g_t(\theta)=U_t(R,\theta)=O(t^5)\) in \(C^1\)
(indeed in every fixed \(C^k\) norm), by Taylor\textquotesingle s
theorem for finitely many smooth functions. Choose a smooth fixed cutoff
\(\chi\) that vanishes for arguments at most \(1/2\) and equals 1 near
1. The function

\[
E_t(\rho,\theta)=\chi(\rho/R(\theta))g_t(\theta)
\tag{4.25}
\]

has trace \(g_t\) and \(H^1(\Omega_{t,b})\) norm \(O(t^5)\), since \(R\)
and its derivatives are uniformly controlled and \(R\) stays bounded
away from zero. Thus \(v_t=U_t-E_t\in H^1_0(\Omega_{t,b})\), and for
every \(\phi\in H^1_0(\Omega_{t,b})\),

\[
\left|\int\nabla v_t\cdot\nabla\phi-k(t)^2\int v_t\phi\right|
\le C|t|^5\|\phi\|_{H^1}.
\tag{4.26}
\]

Let \(\phi_t\) be the positive \(L^2\)-normalized first eigenfunction.
Inclusion between disks of radii \(1\pm C|t|\) gives
\(\lambda_1(\Omega_{t,b})\to j^2\) and a uniform \(H^1\) bound on
\(\phi_t\). Their zero extensions are bounded in \(H^1_0(B_2)\);
compactness gives strong \(L^2\) subsequential convergence. Every limit
vanishes outside \(B\), has norm 1, belongs to \(H^1_0(B)\), and has
energy at most \(j^2\). The Rayleigh principle and positivity identify
it uniquely as the normalized first disk eigenfunction. Hence the whole
family converges to that eigenfunction. Also \(v_t\), extended by zero,
converges in \(L^2\) to \(J_0(j\rho)\) on \(B\) and zero outside. It
follows that \(\int v_t\phi_t\) converges to a positive constant.

Taking \(\phi=\phi_t\) in (4.26) therefore gives

\[
|\lambda_1(\Omega_{t,b})-k(t)^2|=O(t^5).
\tag{4.27}
\]

This proves the eigenvalue line of (4.4), including the remainder.
Together with Sections 4.2 and 4.3 it completes the proof of (4.3).

\Needspace{8\baselineskip}\section{5. The planar perimeter
endpoint}\label{the-planar-perimeter-endpoint}

For \(n=2\), every degree-two harmonic has zero cubic moment, so Section
3.3 does not settle the endpoint. Section 4 already proves the volume
endpoint by a fourth-order calculation. The same paths also settle the
perimeter endpoint.

Write \(s=j_{0,1}^2\), \(q_P=(s+1)/10\), and use \(\Omega_{t,b}\) from
Section 4, with radial function \(R=1+t\cos2\theta+bt^2\cos4\theta\).
Direct expansion of \(\langle\sqrt{R^2+(R')^2}\rangle\) gives

\[
\frac{P(\Omega_{t,b})}{2\pi}
=1+t^2+\left(4b^2-\frac32b-\frac12\right)t^4+O(t^5).
\tag{5.1}
\]

Combine this with the eigenvalue and torsion expansions proved in
Section 4. The quadratic coefficient of
\(\log(G_q(\Omega_{t,b})/G_q(B_1))\) is \((s+1-10q)/2\), which vanishes
at \(q_P\). At that exponent the quartic coefficient is

\[
\begin{aligned}
640(6-s)D_b(s)={}&1264b^2s^2-8960b^2s+9216b^2\\
&+120bs^3-2776bs^2+13760bs-9984b\\
&+35s^4-510s^3+2691s^2-5560s+1824.
\end{aligned}
\tag{5.2}
\]

For \(b=0,1\) the numerator becomes, respectively,

\[
\begin{aligned}
Q_0(s)&=35s^4-510s^3+2691s^2-5560s+1824,\\
Q_1(s)&=35s^4-390s^3+1179s^2-760s+1056.
\end{aligned}
\tag{5.3}
\]

Section 4 gives the exact bracket \(23/4<s<29/5<6\). Setting
\(y=s-23/4\),

\[
\begin{aligned}
Q_0(23/4+y)&=35y^4+295y^3+\frac{6693}{8}y^2+\frac{22659}{16}y+\frac{33003}{256},\\
Q_1(23/4+y)&=35y^4+415y^3+\frac{11157}{8}y^2+\frac{11691}{16}y-\frac{55413}{256}.
\end{aligned}
\tag{5.4}
\]

The first polynomial is positive for \(y\ge0\). The second is increasing
there and is negative even at the upper endpoint, where
\(Q_1(29/5)=-22048/125\). Thus \(D_0(s)>0>D_1(s)\). The remainder is
\(O(t^5)\) by the PDE estimates in Section 4 and
Taylor\textquotesingle s theorem for the perimeter integrand. These
opposite quartic signs prove the planar perimeter saddle. Rescaling by
\(2\pi/P(\Omega_{t,b})\) enforces exact perimeter without changing
\(G_q\). This completes Theorem A in every dimension.

\Needspace{8\baselineskip}\section{6. Local bifurcation and branch
directions}\label{local-bifurcation-and-branch-directions}

The planar branch directions proved in this section are

\[
q_F(t)=q_V+c_Vt^2+O(t^4),\qquad
q_G(t)=q_P+\frac25c_Pt^2+O(t^4),
\qquad c_V>0,\ c_P>0.
\tag{6.1}
\]

\subsection{6.1. Solving the nonlinear
equation}\label{solving-the-nonlinear-equation}

Let \(m_2=\langle H^2\rangle\) and let
\(\Pi f=\langle fH\rangle H/m_2\). Write \(h=tH+w\) with
\(w\in X_\perp\). The analytic implicit-function theorem applied to

\[
(I-\Pi)\Psi(q,tH+w)=0
\tag{6.2}
\]

gives a unique analytic solution \(w=w(q,t)\) near \((q_*,0)\). Since
the ball solves the equation for all \(q\), uniqueness gives
\(w(q,0)=0\). Moreover \(D_h\Psi(q,0)H\) is a multiple of \(H\) for
every \(q\). Differentiating (6.2) at \(t=0\) therefore gives
\(w_t(q,0)=0\), so \(w=O(t^2)\) uniformly for nearby \(q\).

The remaining scalar equation is

\[
\phi(q,t)=\frac{\langle\Psi(q,tH+w(q,t)),H\rangle}{m_2}=0.
\tag{6.3}
\]

Since \(\phi(q,0)=0\), the function

\[
\widetilde\phi(q,t)=\int_0^1\partial_t\phi(q,\tau t)\,d\tau
\tag{6.4}
\]

is analytic and satisfies \(\phi=t\widetilde\phi\). At \((q_*,0)\) it
vanishes, and

\[
\partial_q\widetilde\phi(q_*,0)=2\partial_q e_2(q_*)\ne0.
\tag{6.5}
\]

The scalar implicit-function theorem gives a unique analytic \(q=q(t)\),
and hence \(h(t)=tH+w(q(t),t)\). Together with the ball branch, these
exhaust the solutions in the indicated symmetry class near \((q_*,0)\).
Equation (2.4) proves unrestricted criticality. Since
\(h(t)/t\to H\ne0\) and its mean is zero, the domains are not centered
balls. Their symmetry forces the center of any ball to be the origin, so
they are not translated balls either. Strict convexity follows from
convergence in \(C^2\) to the sphere.

\subsection{\texorpdfstring{6.2. Direction of the branches for
\(n\ge3\)}{6.2. Direction of the branches for n\textbackslash ge3}}\label{direction-of-the-branches-for-nge3}

Let \(C_*=C_V(n,x)\) or \(C_P(n,x)\) from Section 3. On the complement
solution (6.2), define the reduced energy
\(\mathcal E(q,t)=E(q,tH+w(q,t))-E(q,0)\). Because \(L H=0\) at
\(q=q_*\) and \(w=O(t^2)\), its cubic coefficient at the critical
parameter is unchanged by eliminating \(w\). More explicitly the
quadratic form has no \(H\)-\(w\) cross term for any \(q\), by (2.6),
and the \(w\)-\(w\) term has order four. Therefore

\[
\mathcal E(q,t)=e_2(q)m_2t^2+C_*m_3t^3
+O(t^4+|q-q_*|\,|t|^3),\qquad m_3=\langle H^3\rangle.
\tag{6.6}
\]

The range equation and \(\langle H,w_t\rangle=0\) imply
\(\partial_t\mathcal E=m_2\phi\). Differentiating (6.6) and inserting
\(q(t)\) gives

\[
q'(0)=-\frac{3C_*m_3}{2\partial_q e_2\,m_2}.
\tag{6.7}
\]

For the chosen \(H\), Section 3 and the fourth spherical moment give

\[
m_2=\frac{2(n-1)}{n^2(n+2)},\qquad
m_3=\frac{8(n-1)(n-2)}{n^3(n+2)(n+4)}>0.
\tag{6.8}
\]

The two cubic coefficients are nonzero by the classical transcendence
theorem explained in Section 3. Consequently \(q'(0)\ne0\). In explicit
form,

\[
q_F'(0)=\frac{3C_Vm_3}{2(n+2)m_2},\qquad
q_G'(0)=\frac{3(n-1)C_Pm_3}{(n+2)(3n-1)m_2}.
\tag{6.9}
\]

Thus the nonspherical curves cross the threshold rather than only
touching it. Nonvanishing uses a theorem for every rational Bessel
order; it is not extrapolated from a finite-dimensional computation.

\subsection{6.3. The planar fourth-order
reduction}\label{the-planar-fourth-order-reduction}

Now \(H=\cos2\theta\) and \(m_2=1/2\). Rotation through \(\pi/2\)
preserves the reflection-invariant spaces and sends \(H\) to \(-H\).
Uniqueness in the preceding reduction gives \(q(-t)=q(t)\) and relates
\(h(-t)\) to the rotated \(h(t)\). In particular \(q'(0)=0\).

At the critical parameter the quadratic term in \(w(q_*,t)\) contains
only the mode \(\cos4\theta\). To justify the mode selection, the second
derivative of the full rotation-equivariant boundary map is bilinear.
Applied twice to \(\cos2\theta\), its complex Fourier weights can only
be \(0,4,-4\). Projection onto mean-zero functions eliminates weight
zero, and reflection symmetry eliminates the sine term. The inverse
complement linearization preserves the remaining mode. Thus

\[
w(q_*,t)=b_*t^2\cos4\theta+O(t^3).
\tag{6.10}
\]

Let \(C_b(s)\) be the quartic polynomial in \(b\) in Section 4, and
\(D_b(s)\) the one in Section 5, with \(s=j_{0,1}^2\). The fourth-order
coefficient of the reduced energy is respectively \(C_{b_*}\) or
\(D_{b_*}\). An \(O(t^3)\) remainder in (6.10) cannot change this
coefficient: in the quadratic Taylor term it would pair with \(H\),
which is in the Hessian kernel, or enter at order five; the cubic and
higher Taylor terms likewise first change at order five.

The mode-four component of (6.2) at order \(t^2\) is exactly the
derivative of this quartic polynomial with respect to \(b\). For example
\(\partial_b E(q_*,tH+bt^2\cos4\theta)=t^2D_hE[\cos4\theta]\); dividing
its leading order \(t^4\) identifies the equation. Therefore

\[
b_F=\mathop{\rm argmax}_{b\in\mathbb R}C_b(s),\qquad
b_G=\mathop{\rm argmax}_{b\in\mathbb R}D_b(s),
\quad c_V=C_{b_F}(s),\quad c_P=D_{b_G}(s).
\tag{6.11}
\]

These maximizers are unique. The \(b^2\) coefficients are
\(e_4(q_*)\langle\cos^24\theta\rangle<0\) by (2.7). Equivalently they
can be read directly from the displayed polynomials. The strict
positivity proved in Sections 4 and 5 at \(b=0\) gives

\[
c_V\ge C_0(s)>0,\qquad c_P\ge D_0(s)>0.
\tag{6.12}
\]

For a fully explicit specification,

\[
\begin{aligned}
b_F&=-\frac{3s^3-67s^2+320s-192}{4s(11s-64)},\\
b_G&=-\frac{15s^3-347s^2+1720s-1248}{4(79s^2-560s+576)}.
\end{aligned}
\tag{6.13}
\]

The denominators are nonzero by the strictly negative degree-four
Hessian. From (2.6) the planar quadratic terms are \(-2(q-q_V)t^2\) and
\(-5(q-q_P)t^2\). The reduced critical equation consequently reads

\[
-4(q-q_V)t+4c_Vt^3+O(t^5)=0,
\qquad
-10(q-q_P)t+4c_Pt^3+O(t^5)=0
\tag{6.14}
\]

along the corresponding even analytic branch. In writing the remainders
we used \(q(t)-q_*=O(t^2)\) and rotation symmetry. Dividing by \(t\)
gives precisely (6.1). The planar branches lie strictly on the \(q>q_*\)
side for small nonzero \(t\).

\Needspace{8\baselineskip}\section{7. Classification without a symmetry
assumption}\label{classification-without-a-symmetry-assumption}

\subsection{7.1. A slice for translations and
dilations}\label{a-slice-for-translations-and-dilations}

Write \(\langle\cdot\rangle\) for normalized spherical average.
Normalize a near-ball domain by moving its centroid to zero and then
making the mean of its radial function equal to one. These operations
preserve criticality because both functionals are translation and scale
invariant. A sufficiently small \(C^{k,\gamma}\) perturbation of the
sphere remains star-shaped after this centering. The new radial function
is still \(C^{k,\gamma}\) and close to one: transversality of rays to
the nearby boundary and the implicit-function theorem in local sphere
charts give this assertion. No differentiability in a translation
parameter as an operator on the same Hölder space is needed here.

Define

\[
\mathcal X=\{\zeta\in C^{k,\gamma}(S^{n-1}):
\langle\zeta\rangle=0,\ \langle\zeta\omega\rangle=0\}.
\tag{7.1}
\]

For \(\zeta\) small, the centroid condition can be imposed by an
analytic vector \(a=a(\zeta)\in\mathbb R^n\):

\[
R(\zeta)=1+\zeta+a(\zeta)\cdot\omega,
\qquad \langle R(\zeta)^{n+1}\omega\rangle=0.
\tag{7.2}
\]

At \((\zeta,a)=(0,0)\) the derivative of the second expression with
respect to \(a\) is \((n+1)I/n\). Its derivative in
\(\zeta\in\mathcal X\) is zero. The finite-dimensional implicit-function
theorem with Banach parameters yields a unique analytic \(a\), with
\(a(0)=0\) and \(Da(0)=0\). The construction is \(O(n)\)-equivariant. If
\(\zeta\) is centrally even, then \(a(\zeta)=0\) by uniqueness.

Set \(\mathcal E(q,\zeta)=\log F_q(\Omega_{R(\zeta)})\) or
\(\log G_q(\Omega_{R(\zeta)})\). The analytic construction in Section
2.1 applies without a symmetry restriction. Its gradient on the slice is
an analytic map to the corresponding space
\(\mathcal Y=C^{k-r,\gamma}\cap\{\ell=0,1\text{ modes vanish}\}\), where
\(r=1\) for \(F\) and \(r=2\) for \(G\).

To see the gradient assertion explicitly, let \(Z\) be the radial
gradient in (2.2). Differentiating (7.2) expresses \(Da(\zeta)[v]\) as
the inverse of the finite matrix

\[
(n+1)\langle R^n\omega\otimes\omega\rangle
\tag{7.3}
\]

applied to \(-(n+1)\langle R^n v\omega\rangle\). Thus
\(D_\zeta\mathcal E[v]=\langle Z,v+Da[v]\cdot\omega\rangle\) is
represented by \(Z\) plus a finite-rank analytic correction with
\(C^{k,\gamma}\) coefficients. Projection away from degrees zero and one
gives the claimed analytic gradient \(\mathcal G\) with no loss beyond
\(r\) derivatives.

Zero gradient on this slice is equivalent to unrestricted criticality.
Indeed annihilating the tangent space to the two constraints means that
the radial derivative is a linear combination of the derivative of mean
radius and of the centroid moment \(M(\Omega)=\int_\Omega x\,dx\). At a
centered normalized domain, dilation has mean-radius derivative one and
moment derivative \((n+1)M=0\). Scale invariance therefore makes the
mean-radius multiplier zero. Translation by a vector \(v\) has moment
derivative \(|\Omega|v\). Translation invariance then makes every moment
multiplier zero. Hence the unrestricted first variation vanishes.
Conversely unrestricted criticality plainly implies slice criticality.

\subsection{7.2. Reduction to trace-free symmetric
matrices}\label{reduction-to-trace-free-symmetric-matrices}

The linearization of \(\mathcal G\) at the ball has precisely the
harmonic multipliers from (2.6) in Section 2.2, since \(Da(0)=0\) and
the ball is unrestricted critical. At \(q=q_*\) its kernel is the entire
space \(\mathcal H_2\) of degree-two harmonics; every degree
\(\ell\ge3\) has a strictly negative multiplier. The complement
isomorphism proof in Section 2.2 works on all harmonics of degrees at
least three, not just invariant ones: the same Dirichlet-to-Neumann
compact perturbation, index-zero argument and harmonic testing apply.

Identify \(\mathcal H_2\) with the space \(\mathrm{Sym}_0(n)\) of real
trace-free symmetric matrices via

\[
H_A(\omega)=\omega^TA\omega.
\tag{7.4}
\]

Let \(\Pi_2\) be the harmonic projection and solve

\[
(I-\Pi_2)\mathcal G(q,H_A+w)=0,\qquad w\perp\mathcal H_2.
\tag{7.5}
\]

The analytic implicit-function theorem supplies a unique \(w=w(q,A)\),
with \(w(q,0)=0\), \(D_Aw(q,0)=0\), and \(w=O(|A|^2)\). Uniqueness makes
\(w\) equivariant under the action of \(O(n)\). In particular, the
transformation \(-I\) fixes every \(A\), so \(w(q,A)\) is centrally
even. Consequently the centroid correction \(a(H_A+w)\) vanishes
identically on the reduced family. The slice used to remove translations
therefore does not alter any coefficient computed on the even paths in
the preceding calculations.

Define the reduced energy

\[
\mathscr E(q,A)=\mathcal E(q,H_A+w(q,A))-\mathcal E(q,0).
\tag{7.6}
\]

Its critical points correspond exactly to the unrestricted critical
domains near the ball. The implication follows by differentiating (7.6):
the complement gradient vanishes by (7.5), and \(D_Aw\) is orthogonal to
\(\mathcal H_2\). It is \(O(n)\)-invariant under conjugation of \(A\).

The spherical moments are

\[
\langle H_A^2\rangle=\frac{2\operatorname{tr}A^2}{n(n+2)},
\qquad
\langle H_A^3\rangle=\frac{8\operatorname{tr}A^3}{n(n+2)(n+4)}.
\tag{7.7}
\]

They follow by diagonalizing \(A\), expanding the powers, and using the
moments of squared sphere coordinates; the terms containing
\(\operatorname{tr}A\) vanish. The endpoint calculations of Section 3,
valid for every degree-two harmonic, and the kernel property imply

\[
\begin{gathered}
\mathscr E(q,A)=\kappa_2 e_2(q)\operatorname{tr}A^2
+\kappa_3 C_*\operatorname{tr}A^3
+O(|A|^4+|q-q_*||A|^3),\\
\kappa_2=\frac2{n(n+2)},\quad
\kappa_3=\frac8{n(n+2)(n+4)}.
\end{gathered}
\tag{7.8}
\]

Here \(C_*\) is \(C_V(n,j^2/n)\) or \(C_P(n,j^2/n)\). For \(n\ge3\) it
is nonzero by the classical Bessel-zero transcendence result explained
in Section 3. Equation (7.8) is an analytic Taylor expansion, so its
derivatives have the corresponding differentiated orders. The cubic
coefficient is unchanged by the complement correction because the
Hessian vanishes on \(\mathcal H_2\) and \(w=O(|A|^2)\).

\subsection{7.3. Excluding three distinct matrix
eigenvalues}\label{excluding-three-distinct-matrix-eigenvalues}

We give the finite-dimensional argument explicitly; no novel general
invariant-theory theorem is being claimed. Extend the reduced energy to
diagonal matrices of arbitrary trace by setting

\[
f(q,z_1,\ldots,z_n)=\mathscr E(q,\operatorname{diag}(z_i-\bar z)),
\qquad \bar z=\frac1n\sum_i z_i.
\tag{7.9}
\]

It is analytic and symmetric in the \(z_i\). Let
\(g_i=\partial_{z_i}f\). Symmetry implies that \(g_i-g_j\) vanishes on
\(z_i=z_j\), and analytic division by that linear coordinate gives

\[
g_i-g_j=(z_i-z_j)B_{ij}(q,z).
\tag{7.10}
\]

For trace-zero \(z\), equation (7.8) gives

\[
B_{ij}=2\kappa_2e_2(q)+3\kappa_3C_*(z_i+z_j)+R_{ij},
\tag{7.11}
\]

where \(R_{ij}=O(|z|^2+|q-q_*||z|)\). These are analytic remainders, not
just pointwise estimates. The quotient is symmetric in \(i,j\) and
equivariant under permutations. Thus \(R_{ij}-R_{i\ell}\) vanishes
whenever \(z_j=z_\ell\). A second analytic division yields

\[
B_{ij}-B_{i\ell}=(z_j-z_\ell)
\big[3\kappa_3C_*+O(|z|+|q-q_*|)\big].
\tag{7.12}
\]

The remainder order can also be read term by term from the convergent
analytic expansion: the higher energy terms have degree at least four in
\(z\), or have a factor \(q-q_*\) and degree at least three; one
differentiation and two divisions lower these degrees by three. If
desired, each analytic division is justified directly by the integral
formula for the difference of an analytic function at two values of one
coordinate.

At a critical trace-free matrix, \(g_i=0\) for every \(i\). Indeed the
derivative on the trace-free diagonal subspace vanishes and
\(\sum_i g_i=0\) by construction. If three eigenvalues
\(z_i,z_j,z_\ell\) were distinct, (7.10) would give
\(B_{ij}=B_{i\ell}=0\). Equation (7.12) would contradict \(C_*\ne0\)
once \((q,z)\) is sufficiently close to \((q_*,0)\). Hence every nonzero
critical matrix has at most two distinct eigenvalues.

This proves the necessary exclusion for \(n\ge3\). For \(n=2\) every
trace-free symmetric matrix already has at most two eigenvalues, so no
cubic argument is required. A nonzero trace-free matrix cannot have only
one eigenvalue in any dimension.

\subsection{7.4. The finite list of branches and their
slopes}\label{the-finite-list-of-branches-and-their-slopes}

Every trace-free matrix with exactly two eigenvalues is orthogonally
conjugate to

\[
A=t\left(\operatorname{diag}(I_p,0)-\frac pn I_n\right)
\tag{7.13}
\]

for some \(1\le p\le n-1\) and \(t\ne0\). Exchanging the two blocks
replaces \((p,t)\) by \((n-p,-t)\), so the list may be restricted to
\(p\le\lfloor n/2\rfloor\). Its isotropy group contains
\(K_p=O(p)\times O(n-p)\). Equivariance of the unique complement
solution makes the full radial graph \(K_p\)-invariant.

The fixed subspace of \(\mathcal H_2\) under \(K_p\) is exactly the line
spanned by \(H_p\). Moreover the reduced gradient at \(tH_p\) belongs to
this fixed line. Thus unrestricted criticality there is equivalent to
its scalar equation along \(t\). The construction of Sections 2.2 and
6.1 applies with \(K\) replaced by \(K_p\); the kernel remains
one-dimensional and \(\partial_q e_2\ne0\). It gives a unique analytic
curve \(q_p(t)\) and no other nontrivial solution near the origin in
this class. Because every critical matrix has the form (7.13), this list
exhausts all nearby critical domains.

For the moments one can use the beta distribution of
\(\sum_{i=1}^p\omega_i^2\), whose parameters are \(p/2,(n-p)/2\). It
gives

\[
m_{2,p}=\frac{2p(n-p)}{n^2(n+2)},\qquad
m_{3,p}=\frac{8p(n-p)(n-2p)}{n^3(n+2)(n+4)}.
\tag{7.14}
\]

The slope formula from Section 6 is therefore

\[
q_p'(0)=-\frac{3C_*m_{3,p}}{2\partial_qe_2\,m_{2,p}},
\tag{7.15}
\]

which is nonzero for \(p<n/2\). When \(p=n/2\), an orthogonal
transformation exchanges the equal-sized coordinate blocks and sends
\(H_p\) to \(-H_p\). Uniqueness makes \(q_p(t)\) even and identifies the
two signs of \(t\) geometrically. Section 8 proves

\[
q_p(t)=q_*+\eta_nt^2+O(t^4),\qquad \eta_n\ne0,
\tag{7.16}
\]

by an explicit fourth-order calculation. For \(n=2\),
\(H_1=\tfrac12\cos2\theta\), so the planar amplitude in Section 6 equals
half the parameter used here. Its positive coefficients therefore yield
\(\eta_2>0\) for both functionals.

\subsection{7.5. Consequences and precise
limits}\label{consequences-and-precise-limits}

Using the balanced calculation in Section 8, the ball is an isolated
critical domain at \(q=q_*\) modulo similarities, even though it is a
nonlinear saddle at that parameter. For unbalanced branches this follows
from (7.15); for the balanced branch it follows from (7.16). Finitely
many branches allow one common sufficiently small neighborhood.

In the plane, for every fixed \(q>q_*\) sufficiently close to \(q_*\),
there is exactly one nearby nonball critical domain up to similarities.
For \(q\le q_*\) sufficiently close there is none. Indeed the even
analytic curve has the form \(q(t)=Q(t^2)\) with \(Q'(0)>0\), so it is
invertible as a function of \(t^2\ge0\); the two signs represent
rotations of the same shape.

In odd dimensions there are \((n-1)/2\) nearby nonball similarity
classes at each sufficiently close \(q\ne q_*\). In even dimensions at
least four, the \(n/2-1\) unbalanced classes occur on both sides, and
the balanced class adds one on the side determined by the nonzero
coefficient \(\eta_n\). No universal sign of that balanced coefficient
in dimensions at least four is asserted here. Different \(p<n/2\) give
distinct eigenvalue multiplicities in their degree-two projections and
hence cannot be similar after normalization.

\Needspace{8\baselineskip}\section{8. Balanced fourth-order
expansion}\label{balanced-fourth-order-expansion}

\subsection{8.1. The balanced harmonic}\label{the-balanced-harmonic}

Let \(n\ge2\) be even, \(s=j_{n/2-1,1}^2\), and set

\[
\begin{gathered}
h=\sum_{i=1}^{n/2}\omega_i^2-\frac12,\qquad Q=h^2-m_2,\\
m_2=\langle h^2\rangle=\frac1{2(n+2)},\qquad
m_4=\langle h^4\rangle=\frac3{4(n+2)(n+6)},\\
N=\langle Q^2\rangle=\frac{n}{2(n+2)^2(n+6)}.
\end{gathered}
\tag{8.1}
\]

The beta distribution with parameters \(n/4,n/4\) proves these moments.
All odd moments vanish. Direct differentiation gives

\[
-\Delta_Sh=2nh,\qquad |\nabla_Sh|^2=1-4h^2,\qquad
-\Delta_SQ=4(n+2)Q.
\tag{8.2}
\]

Thus \(h\) and \(Q\) have harmonic degrees two and four. Consider

\[
R=1+th+bt^2Q.
\tag{8.3}
\]

Block exchange sends \(t\) to \(-t\), so the scalar functionals are
even. We compute the critical quartic coefficients \(K_V(b)\) and
\(K_P(b)\) of the two logarithmic functionals. Their quadratic
coefficients in \(b\) are negative. Their maxima \(c_V(n),c_P(n)\) are
nonzero for every even \(n\). Hence the balanced critical curves satisfy

\[
q(t)=q_*+\eta_nt^2+O(t^4),\qquad
\eta_n=-\frac{2c_*(n)}{\partial_qe_2\,m_2}\ne0.
\tag{8.4}
\]

Only nonvanishing is proved for all even dimensions; the planar sign is
positive by Section 6.

\subsection{8.2. Eigenvalue jet and its analytic
remainder}\label{eigenvalue-jet-and-its-analytic-remainder}

Let \(f_0\) be the regular radial ground state at eigenvalue \(s\),
normalized by \(f_0(1)=0,f_0'(1)=1\). For \(\ell>0\), let \(f_\ell\) be
the regular degree-\(\ell\) radial Helmholtz solution, normalized by
\(f_\ell(1)=1\). Its boundary value is nonzero before normalization
because the first eigenspace is radial. Write

\[
d_2=\frac sn-n,\qquad
d_4=-n-2+\frac{s[n(n+2)-s]}{(n+2)[n(n+4)-2s]}.
\tag{8.5}
\]

These equal \(f_2'(1),f_4'(1)\) by the Bessel recurrence. Denote higher
derivatives at one by \(f_{\ell r}\). The radial ODE gives

\[
\begin{aligned}
f_{02}&=-(n-1),\qquad f_{03}=n(n-1)-s,\\
f_{04}&=(n-1)[2s-n(n+1)],\\
f_{22}&=-(n-1)d_2-s+2n,\\
f_{23}&=-(n-1)f_{22}+(3n-1-s)d_2-4n,\\
f_{42}&=-(n-1)d_4-s+4(n+2).
\end{aligned}
\tag{8.6}
\]

Set

\[
a=d_2+\frac{n-1}{2},\quad A=am_2,\quad a_4=a-b,
\tag{8.7}
\]

\[
c=-A(1+f_{02}-d_2)
-[b(f_{02}-d_2)+a_4d_4]\frac{N}{m_2}
-\left(\frac{f_{03}}6-\frac{f_{22}}2\right)\frac{m_4}{m_2},
\tag{8.8}
\]

and define \(B\) by

\[
\begin{aligned}
-B={}&\frac{f_{02}}2(b^2N+A^2+2Am_2)
+\frac{f_{03}}2(bN+Am_2)+\frac{f_{04}}{24}m_4\\
&-d_2Am_2-f_{22}(bN+Am_2)-\frac{f_{23}}6m_4
+a_4\left(d_4b+\frac{f_{42}}2\right)N+cd_2m_2.
\end{aligned}
\tag{8.9}
\]

Then

\[
\log\frac{\lambda_1(\Omega_t)}s
=2At^2+(2B-A^2)t^4+O(t^6).
\tag{8.10}
\]

To prove this, put \(\rho=(1+At^2+Bt^4)r\) and construct a finite sum
beginning with

\[
U=f_0(\rho)-tf_2(\rho)h+t^2a_4f_4(\rho)Q
+t^3[cf_2(\rho)h+f_6(\rho)H_6]+t^4U_4.
\tag{8.11}
\]

Here \(H_6\) is a degree-six harmonic and \(U_4\) is a sum of
degree-four and degree-eight regular Helmholtz modes. At the moving
boundary,

\[
\rho-1=th+t^2(bQ+A)+At^3h+t^4(B+AbQ)+O(t^5).
\tag{8.12}
\]

The constant and degree-four equations at order two give (8.7). At order
three only degrees two and six occur. Projecting onto degree two, using
\(\Pi_2(hQ)=(N/m_2)h\) and \(\Pi_2(h^3)=(m_4/m_2)h\), gives (8.8); the
remaining equation determines \(H_6\).

At order four, take the boundary average. The three terms in the first
line of (8.9) come from \(f_0\). The first three terms of the second
line come from \(-tf_2h\). The last two come from the degree-four and
corrected degree-two terms. The degree-six term pairs with \(h\) to
zero. Hence the mean equation gives (8.9). The remaining modes at this
order have degrees four and eight and are cancelled by \(U_4\).

Every term solves the Helmholtz equation exactly at eigenvalue
\(s(1+At^2+Bt^4)^2\). The remaining Dirichlet data are \(O(t^5)\) in a
fixed smooth norm. Subtract an \(H^1\) boundary lift of that order and
test against the normalized first eigenfunction on \(\Omega_t\). The
pairing has a nonzero limit, as in Section 4, so the eigenvalue error is
\(O(t^5)\). Analytic dependence from Section 2.1 and block symmetry make
the eigenvalue an even function of \(t\), improving the scalar remainder
to \(O(t^6)\). Taking logarithms proves (8.10). This argument justifies
the jet analytically, not just formally.

\subsection{8.3. Torsion and geometric
jets}\label{torsion-and-geometric-jets}

Write the normalized torsion, volume and perimeter as
\(1+X_2t^2+X_4t^4+O(t^6)\). Their quadratic coefficients are

\[
T_2=\frac{(n+2)(n-3)}2m_2,\quad
V_2=\frac{n(n-1)}2m_2,\quad
P_2=\frac{n^2-n+2}2m_2.
\tag{8.13}
\]

Their quartic coefficients are

\[
\begin{aligned}
T_4=(n+2)\bigg\{&
\left[\frac{n-7}2b^2+\frac{n^2-7n+16}2b
-\frac{3(n-1)(n-4)}4\right]N\\
&+\frac{48+n^3+10n^2-27n}{24}m_4-\frac{3n(n-1)}4m_2^2\bigg\},\\
V_4={}&\frac{n(n-1)}2b^2N+\frac{n(n-1)(n-2)}2bN
+\frac{n(n-1)(n-2)(n-3)}{24}m_4,\\
P_4={}&\left[\frac{(n-1)(n-2)}2+2(n+2)\right]b^2N\\
&+\left[\frac{(n-1)(n-2)(n-3)}2+2(n-3)(n+1)\right]bN\\
&+\left[\frac{(n-1)(n-2)(n-3)(n-4)}{24}
+\frac{n(n-3)(n-4)}6\right]m_4
-\frac{1-8m_2+16m_4}{8}.
\end{aligned}
\tag{8.14}
\]

Here is a derivation of the torsion formula. The first corrections are

\[
w_0=\frac{1-r^2}{2n},\quad w_1=\frac{r^2h}{n},\quad
w_2=\frac{(b-3/2)r^4Q-(3/2)m_2}{n},
\quad
w_3|_S=\frac{(6-5b)hQ-h^3}{n}.
\tag{8.15}
\]

They cancel the moving boundary through order three. At order four,
harmonic orthogonality gives
\(\langle h\partial_rw_3\rangle=2[(6-5b)N-m_4]/n\), and the mean
boundary correction becomes

\[
\langle w_4|_S\rangle
=\frac{(-7b^2/2+8b-3)N+2m_4}{n}.
\tag{8.16}
\]

The fourth-order signed shell integral, divided by \(|S^{n-1}|\), is

\[
\frac1n\left\{\left[\frac{b^2}2+\frac{n-7}2b+\frac{15-3n}4\right]N
+\frac{n^2+10n-27}{24}m_4-\frac{3(n-1)}4m_2^2\right\}.
\tag{8.17}
\]

For verification, the four contributions before collection, multiplied
by \(n/|S|\), are \(-b^2N/2-(2n-1)bN/2-(n-1)^2m_4/8\) from \(w_0\);
\((n+1)bN+n(n+1)m_4/6\) from \(tw_1\);
\((b-3/2)bN+(b-3/2)(n+3)N/2-3(n-1)m_2^2/4\) from \(t^2w_2\); and
\((6-5b)N-m_4\) from \(t^3w_3\). Adding the ball integral
\(|S|\langle w_4|_S\rangle/n\) and dividing by \(T(B)=|S|/[n^2(n+2)]\)
gives \(T_4\).

The corrections are finite harmonic polynomials. Their sum has Laplacian
exactly \(-1\) and uncancelled boundary value \(O(t^5)\). The maximum
principle controls the true torsion error. Analyticity and evenness
improve the scalar integral remainder to \(O(t^6)\).

Volume follows by expanding \(\langle R^n\rangle\). For perimeter expand
\(\langle R^{n-1}\sqrt{1+|\nabla R|^2/R^2}\rangle\) and use

\[
\begin{gathered}
\langle|\nabla Q|^2\rangle=4(n+2)N,\quad
\langle h\nabla h\cdot\nabla Q\rangle=2(n+2)N,\quad
\langle Q|\nabla h|^2\rangle=-4N,\\
\langle h^2|\nabla h|^2\rangle=\frac{2n}3m_4,\qquad
\langle|\nabla h|^4\rangle=1-8m_2+16m_4.
\end{gathered}
\tag{8.18}
\]

These identities follow from (8.2) and integration by parts. In
particular the final negative term of \(P_4\) comes from the quartic
term in the square-root expansion.

\subsection{8.4. Complement elimination and an explicit rational
coefficient}\label{complement-elimination-and-an-explicit-rational-coefficient}

Use \(q_V,q_P\) from (1.2), \(\alpha_V=((n+2)q_V-2)/n\) and
\(\beta_P=((n+2)q_P-2)/(n-1)\). The two critical quartics are explicitly
specified by

\[
\begin{aligned}
K_V(b)&=2B-A^2+q_V(T_4-T_2^2/2)-\alpha_V(V_4-V_2^2/2),\\
K_P(b)&=2B-A^2+q_P(T_4-T_2^2/2)-\beta_P(P_4-P_2^2/2).
\end{aligned}
\tag{8.19}
\]

The coefficient of \(b^2\) is \(e_4(q_*)N<0\) by the established
degree-four Hessian. In this balanced symmetry class, the quadratic
complement term must be a multiple of \(Q\): products of two degree-two
harmonics have only degrees zero, two and four; block exchange excludes
degree two here, and the mean-radius normalization excludes zero. The
complement equation selects the unique maximum of \(K_*\).

Writing \(K_*=u_*b^2+v_*b+w_*\), the reduced coefficient is therefore

\[
b_*=-\frac{v_*}{2u_*},\qquad c_*(n)=w_*-\frac{v_*^2}{4u_*}.
\tag{8.20}
\]

Higher complement terms do not affect this quartic coefficient because
the first velocity belongs to the Hessian kernel. This is the same
argument as in Section 6.

To prove nonvanishing without a dimension-by-dimension numerical test,
define

\[
\begin{aligned}
D_V&=-2n^3-16n^2-24n+(5n+12)s,\\
D_P&=(15n^2+43n+12)s^2
-(6n^4+52n^3+116n^2+72n)s+12n^4+72n^3+96n^2.
\end{aligned}
\tag{8.21}
\]

Rational simplification of (8.5)-(8.20) gives

\[
c_V(n)=\frac{N_V(n,s)}{4n^3s(n+2)^2(n+6)D_V},\qquad
c_P(n)=\frac{N_P(n,s)}{4n^3(n+2)^2(n+6)(3n-1)D_P}.
\tag{8.22}
\]

Here \(N_*=\sum_{r=0}^5 a_{*,r}s^r\), with coefficients

\[
\begin{aligned}
a_{V,0}&=-16n^6(n+2)(n+4),\\
a_{V,1}&=16n^5(n+2)(3n+16),\\
a_{V,2}&=-2n^3(n^4+34n^3+208n^2+136n-288),\\
a_{V,3}&=n^2(17n^3+94n^2-512n-1152),\\
a_{V,4}&=-2n(2n^3+3n^2-214n-360),\\
a_{V,5}&=8(n^2-10n-18),
\end{aligned}
\tag{8.23}
\]

and

\[
\begin{aligned}
a_{P,0}&=-4n^5(n+2)(n+4)(9n^3-6n^2+49n-36),\\
a_{P,1}&=4n^4(n+2)(66n^4+269n^3-565n^2-548n+108),\\
a_{P,2}&=-n^3(18n^6+471n^5+1980n^4-4895n^3-19470n^2-8596n+4608),\\
a_{P,3}&=n^2(117n^5+168n^4-8229n^3-18886n^2-4958n+3960),\\
a_{P,4}&=-2n(3n-1)(6n^4-5n^3-753n^2-1622n-648),\\
a_{P,5}&=8(3n-1)(3n^3-31n^2-68n-18).
\end{aligned}
\tag{8.24}
\]

These algebra identities follow by substitution in (8.19) and completing
the square as in (8.20).

\subsection{8.5. Nonvanishing and the critical
curve}\label{nonvanishing-and-the-critical-curve}

For each integer \(n\ge2\), both polynomials \(N_V(n,s)\) and
\(N_P(n,s)\) are nonzero rational polynomials in \(s\). Their constant
terms in (8.23)-(8.24) suffice: the first is negative, and in the second
\(9n^3-6n^2+49n-36=n^2(9n-6)+(49n-36)>0\).

The classical Siegel theorem, whose applicable statement is documented
in Section 3, makes the nonzero Bessel zero \(j_{n/2-1,1}\)
transcendental. Its square \(s\) is also transcendental. Thus neither
numerator in (8.22) can vanish. Neither denominator can vanish: its
nonconstant factors are nonzero rational polynomials in \(s\), since
\(5n+12>0\) and \(15n^2+43n+12>0\). The intermediate denominator in
(8.5) is also nonzero, consistently with spectral simplicity.

Consequently \(c_V(n)c_P(n)\ne0\). Differentiating the reduced energy
\(e_2(q)m_2t^2+c_*t^4+\cdots\) and using \(\partial_qe_2\ne0\) proves
(8.4). This completes the balanced nondegeneracy step in Theorem B and
Corollary C.

For \(n=2\), \(h=\cos2\theta/2\), \(Q=\cos4\theta/8\). Substituting
\(t_{\rm planar}=t/2\) and \(b_{\rm planar}=b/2\) recovers the earlier
planar eigenvalue, torsion and perimeter jets exactly. The reduced
quartic is one sixteenth of the earlier one, so it is positive. The sign
in every higher even dimension is left open.

\Needspace{8\baselineskip}\section{9. Local versus global
maximality}\label{local-versus-global-maximality}

For either \(F_q\) or \(G_q\) of Section 3, let \(q_*\) be its upper
local stability threshold. In every dimension \(n\ge2\), there is
\(\varepsilon_n>0\) such that for every

\[
q_*<q<q_*+\varepsilon_n
\tag{9.1}
\]

the ball is a strict local maximizer, modulo its Euclidean and scaling
symmetries, but is not a global maximizer even among smooth strictly
convex domains. The local topology here is the \(C^{2,\gamma}\) topology
used in the primary source. The constant can differ for the two
functionals. No numerical size is asserted.

\textbf{Proof.} The endpoint theorem proved in Sections 3-5 supplies a
fixed smooth strictly convex domain \(\Omega_*\) with

\[
\log\frac{F_{q_*}(\Omega_*)}{F_{q_*}(B)}>0
\quad\text{or}\quad
\log\frac{G_{q_*}(\Omega_*)}{G_{q_*}(B)}>0,
\tag{9.2}
\]

as appropriate. Choose a sufficiently small but nonzero deformation from
that proof, then keep it fixed. For this domain the logarithmic ratio is
an affine, hence continuous, function of \(q\). Its strict positivity
persists for all \(q\) in an interval containing \(q_*\). This proves
failure of global optimality in (9.1). Local strict maximality for every
\(q>q_*\) is already Theorem 1.6 or Theorem 6.3 of
Amato-Nitsch-Trombetti-Villone,
\href{https://doi.org/10.1007/s11118-026-10328-2}{final article}.
Combining that established result with (9.2) proves the assertion. This
argument does not assert a neighborhood uniform as \(q\downarrow q_*\).
The competitor is fixed, and need not belong to the shrinking
neighborhood on which local maximality holds. \(\square\)

\raggedright\Needspace{8\baselineskip}\section{References}\label{references}

\begin{enumerate}
\def\labelenumi{\arabic{enumi}.}
\tightlist
\item
  V. Amato, C. Nitsch, C. Trombetti and F. Villone, \emph{On Some
  Functionals Involving Torsional Rigidity, Principal Eigenvalue and
  Perimeter}, Potential Analysis 65, article 25 (2026).
  \href{https://doi.org/10.1007/s11118-026-10328-2}{DOI}.
\item
  A. Celentano, C. Nitsch and C. Trombetti, \emph{On a Serrin Type
  Overdetermined Problem}, Milan Journal of Mathematics 92 (2024),
  579-590. \href{https://doi.org/10.1007/s00032-024-00405-9}{DOI}.
\item
  M. G. Crandall and P. H. Rabinowitz, \emph{Bifurcation from simple
  eigenvalues}, Journal of Functional Analysis 8 (1971), 321-340.
  \href{https://doi.org/10.1016/0022-1236(71)90015-2}{DOI}.
\item
  L. Lorch and M. E. Muldoon, \emph{Transcendentality of zeros of higher
  derivatives of functions involving Bessel functions}, International
  Journal of Mathematics and Mathematical Sciences 18 (1995), 551-560.
  \href{https://doi.org/10.1155/S0161171295000706}{DOI}.
\item
  D. Gilbarg and N. S. Trudinger, \emph{Elliptic Partial Differential
  Equations of Second Order}, second edition, Classics in Mathematics,
  Springer (2001).
  \href{https://doi.org/10.1007/978-3-642-61798-0}{Publisher}. Classical
  Schauder theory and elliptic solvability are used as established
  background.
\item
  A. L. Masiello, G. Paoli and F. Salerno, \emph{On some functionals for
  which the ball is a saddle shape}, arXiv:2609.17083v1 (September 15,
  2026). \href{https://arxiv.org/abs/2609.17083}{Primary record}.
\item
  P. Chossat, R. Lauterbach and I. Melbourne, \emph{Steady-state
  bifurcation with O(3)-symmetry}, Archive for Rational Mechanics and
  Analysis 113 (1991), 313-376.
  \href{https://doi.org/10.1007/BF00374697}{DOI}.
\end{enumerate}

\end{document}
